\section{Enquadramentos e objetivos}
\label{sec:enq-e-obj}
%
No âmbito da Unidade Curricular de Sistemas Autónomos Inteligentes do curso de Mestrado em Engenharia Eletrónica Industrial e Computadores da Universidade do Minho, desenvolveu-se um robô autónomo capaz de detetar obstáculos e afastar-se deles, à medida que se aproxima do alvo pretendido. Pretendeu-se desta forma conciliar os conhecimentos adquiridos nas diferentes unidades curriculares da especialização de Controlo, Automação e Robótica e utilizá-los no desenvolvimento do projeto. 

O objetivo concreto deste projeto passou, então, por desenvolver um robô desde a sua raiz, criando e/ou escolhendo todas as peças constituintes deste. Este robô tem de ser capaz de se movimentar fazendo uma trajetória completamente autónoma, tendo apenas como conhecimento inicial as coordenadas do alvo final. Durante este trajeto até ao alvo, o robô tem de ser capaz de agir conforme as situações a que poderá vir a ser sujeito, tais como a presença de obstáculos na sua trajetória ou até mesmo alteração em tempo real da posição do alvo.  

Tendo isto em conta, para o desenvolvimento do robô é possível identificar os seguintes tópicos:

\begin{itemize}
\item Esboço e definição de medidas das várias peças do robô;
\item Escolha e pesquisa das peças que constituem o robô;
\item Desenvolvimento do modelo do robô com uso da ferramenta CAD;
    \begin{itemize}
        \item Criação do modelo físico da plataforma em CoppeliaSim; 
        \item Adaptação dos scripts Matlab para controlo do novo robô;
    \end{itemize}
\item Navegação autónoma:
    \begin{itemize}
        \item -Implementação de algoritmos para a navegação autónoma baseados em sistemas dinâmicos não-lineares;
        \item Desenvolvimento de um programa em Matlab e comunicação entre Matlab e CoppeliaSim para o controlo do robô; 

 
    \end{itemize}
\end{itemize}